\clearpage
\section{Ein historischer Abriss}
\label{bin:history}

Die Geschichte des Pascalschen Dreiecks beginnt lange vor Pascal.
Dabei sind der Ursprung eines Verfahrens, seine erhaltene Überlieferung
und seine spätere systematische Darstellung zu unterscheiden.

\textbf{Vor Pascal.}
Yang Huis 1261 abgeschlossener Kommentar zu den \emph{Neun Kapiteln}
überliefert einen Abschnitt über das Zahlendreieck und ältere
Überlegungen Jia Xians aus dem 11.~Jahrhundert. Jia Xians einschlägiger
Kommentar ist nicht unabhängig erhalten. Die Überlieferung durch
spätere Autoren erlaubt daher keine einfache Gleichsetzung von
Verfasser und erstem Entdecker.\footnote{Zum erhaltenen Dreiecksabschnitt siehe Lam Lay Yong;
zur Überlieferung der Kommentare siehe Chemla, besonders S.~18--19.
Die Nachweise stehen am Ende dieses Abschnitts.}

\textbf{Pascal: aus einem Zahlenschema wird eine Theorie.}
Blaise Pascals \emph{Traité du triangle arithmétique} erschien 1665,
drei Jahre nach seinem Tod. Der Text entwickelt die Bildung der
Einträge und ihre Summationsgesetze und verbindet das Dreieck mit der
Berechnung von Kombinationen. Die Zerlegung aller Auswahlen danach,
ob ein festgehaltenes Element enthalten ist, erklärt die Additionsregel.

Pascal beschreibt auch, wie eine Aussage für das erste Dreieck gezeigt
und vom jeweiligen auf das folgende übertragen wird. Unsere Darstellung
verwendet diese Gedanken in heutiger Sprache und mit der Schreibweise
$\binom nk$.

\textbf{Newton: der Schritt zu unendlichen Reihen.}
Im Zusammenhang mit seiner Lektüre von John Wallis im Winter 1664/65
untersuchte Newton Potenzentwicklungen mit negativen und gebrochenen
Exponenten. Sein
mathematisches Notizbuch bewahrt die Rechnungen. Für natürliche
Exponenten endet die Binomialentwicklung; bei anderen Exponenten
kann eine unendliche Reihe entstehen, deren Konvergenz zu prüfen ist.
Diese Erweiterung führt über den hier bewiesenen endlichen Lehrsatz hinaus.

\textbf{Vandermonde: regelmäßige Produkte.}
Vandermonde untersucht Produkte wie $p(p-1)(p-2)\cdots$ und sucht
dafür Rechenregeln, später auch für gebrochene Ordnungen. Aus einer
Regel für die Änderung bei $p\mapsto p+1$ gewinnt er die Fälle
$p+2$, $p+3$ und die allgemeine Additionsformel. Seine Motivation
ist die Erweiterung der Potenzrechnung mit Anwendungen auf geometrische
Größen. Die Wege in dieser Lesefassung sind eine heutige Erklärung;
sie werden nicht als sein historischer Entdeckungsweg ausgegeben.

\subsection*{Originaltexte und Quellen zur Überlieferung}
\begingroup
\small
\begin{enumerate}
\item Lam Lay Yong: \emph{On the existing fragments of Yang Hui's
Hsiang Chieh Suan Fa}. Archive for History of Exact Sciences~6 (1969),
S.~82--88. Quellenstudie zu den erhaltenen Fragmenten, einschließlich
des Dreiecksabschnitts.
\href{https://doi.org/10.1007/BF00327264}{Verlagsnachweis der Studie}.

\item Karine Chemla: \emph{Reading instructions of the past,
classifying them and reclassifying them}. BJHS Themes~5 (2020),
S.~15--37. Zur Überlieferung von Jia Xians und Yang Huis Kommentaren.
\href{https://doi.org/10.1017/bjt.2020.2}{Offene wissenschaftliche Quellenstudie}.

\item Blaise Pascal: \emph{Traité du triangle arithmétique}.
Paris: Guillaume Desprez, 1665. Hier besonders der Abschnitt über
Kombinationen, Lemma~4 und Proposition~1.
\href{https://gallica.bnf.fr/ark:/12148/btv1b86262012}{Originaldigitalisat der BnF};
\href{https://fr.wikisource.org/wiki/Trait%C3%A9_du_triangle_arithm%C3%A9tique}{Transkription};
\href{https://ccfr.bnf.fr/portailccfr/ark:/16871/00114268519}{Bibliotheksnachweis der Ausgabe}.

\item Isaac Newton: \emph{Mathematical Notebook}.
Cambridge University Library, MS Add.~4000, fol.~14v und folgende.
\href{https://newtonproject.ox.ac.uk/view/texts/diplomatic/NATP00128}{Originaltext in wissenschaftlicher Transkription};
\href{https://newtonproject.ox.ac.uk/texts/introduction}{Einordnung durch das Newton Project}.
\item Alexandre-Théophile Vandermonde: \emph{Abhandlung über irrationale
Grössen verschiedener Ordnung nebst einer Anwendung auf den Kreis}.
In: \emph{Abhandlungen aus der reinen Mathematik}, deutsche Ausgabe von
Carl Itzigsohn, Berlin 1888, S.~67--81, besonders §~VI, S.~71.
Französisches Original: Akademieband 1772, Teil~I, gedruckt 1775,
S.~489--498.
\href{https://www.e-rara.ch/download/pdf/30678740.pdf\#page=6}{Deutsche Übersetzung, entscheidende Seite}.
\end{enumerate}
\endgroup
